\begin{comment}
\begin{frame}{Recherche d'information}

  3 cours : 
  \begin{itemize}
      \setlength{\itemsep}{6pt}
    \item 17 octobre : introduction générale à la RI
    \item 24 octobre : indexation
    \item 31 octobre : recherche avec classement
  \end{itemize}
  \bigskip
  Travaux pratiques :
  \begin{itemize}
      \setlength{\itemsep}{6pt}
    \item 2-3 novembre : MongoDB
    \item 9-10 novembre : MongoDB MapReduce
    \item {\color{red}7-8 décembre : ElasticSearch (1)}
    \item {\color{red}14-15 décembre : ElasticSearch (2)}
    \item 11-12 janvier 2018 : calcul distribué (1)
    \item 18-19 janvier 2018 : calcul distribué (2)
  \end{itemize}
\end{frame}
\end{comment}

\subsection{Exemple de Recherche d'information}
\begin{frame}{\textit{Information Retrieval}, définition}

  \vfill
  \begin{defblock}{Une définition}
    La \alert{Recherche d'Information} (\textit{Information Retrieval}, IR) consiste à
    trouver des \alert{documents} peu ou faiblement structurés, dans une grande \alert{collection},
    en fonction d'un \alert{besoin d'information}.
  \end{defblock}

  \vfill
  \begin{itemize} 
      \setlength{\itemsep}{8pt}
    \item Recherche sur le Web. Utilisée quotidiennement par des milliards d'utilisateurs.
    \item Recherche dans votre boîte mail.
    \item Recherche sur votre ordinateur (\textit{Spotlight}).
    \item Recherche dans une base documentaire, publique ou privée.
  \end{itemize}
  \vfill
\end{frame}

\begin{frame}{\textit{Information Retrieval}, définition}
  \begin{itemize}
    \item Recherche plein texte : on cherche à examiner tous les mots de chaque
      document enregistré et à essayer de les faire correspondre à ceux fournis
      par l'utilisateur
    \item À bien distinguer d'une recherche type ``base de données''\,: requête
      structurée, données structurées, réponse \guill{exacte}.
  \end{itemize}
  \vfill
  
\begin{blocgauche}
  Contrairement à une base de données (SQL), le résultat dépend de
  \alert{l'interprétation} d'un besoin. On ne peut jamais dire qu'un résultat est
  totalement exact (ou totalement faux).
\end{blocgauche}
\end{frame}

\begin{comment}
\begin{frame}{Un exemple de besoin}
  \begin{center}
    \includegraphics[height=7cm]{avocat.jpg}
  \end{center}
\end{frame}

\begin{frame}[plain]{Un moteur de recherche bien connu}
  \begin{center}
    \includegraphics[height=9cm]{serp-googleavocat.png}
  \end{center}
\end{frame}

\begin{frame}{Fonctionnalités avancées}
  \begin{center}
    \includegraphics[height=7cm]{serp-wiki-avocat-autocomplete.png}
  \end{center}
\end{frame}

\begin{frame}{Wikipédia : recherche avancée}
  \begin{center}
    \includegraphics[width=10cm]{serp-wiki-avocat-facettes.png}
  \end{center}
\end{frame}

\begin{frame}{ECommerce : recherche à facettes}
  \begin{center}
    \includegraphics[width=10cm]{serp-houra-avocat-facettes.png}
  \end{center}
\end{frame}

\end{comment}
\subsection{Un besoin, une réponse}

\begin{frame}[plain]{Évaluation d'un moteur de recherche}
  \begin{center}
    \only<1>{\includegraphics[height=8cm]{serp-googleavocat.png}}
    \only<2>{\includegraphics[height=8cm]{serp-googleavocat-evaluation.png}}
  \end{center}
\end{frame}

\begin{frame}{Évaluation d'un moteur de recherche}
  \begin{defblock}{Deux notions importantes}
  \begin{itemize}
    \item \alert{Faux positifs} : ce sont les documents \alert{non pertinents} inclus dans le résultat;
      ils ont été sélectionnés à tort.
    \item \alert{Faux négatifs} : ce sont les documents \alert{pertinents} qui \alert{ne sont pas}
      inclus dans le résultat. 
  \end{itemize}
\end{defblock}

  \vfill
  \begin{itemize}
    \item 
      La recherche plein texte est susceptible de récupérer beaucoup de faux positifs. 
    \item La récupération de documents non pertinents est souvent
      provoquée par l'ambiguïté inhérente au langage naturel ; 
  \end{itemize}
\vfill

\begin{blocgauche}
En général, chercher à réduire les faux positifs entraîne l'augmentation des
  faux négatifs, et réciproquement. 
\end{blocgauche}

\end{frame}

\begin{comment}
\begin{frame}{Évaluation de la pertinence : précision et rappel}

  La \alert{précision} mesure la fraction des vrais positifs dans le résultat $r$.

  \medskip

  Si on note $t_p(r)$ et $f_p(r)$ le nombre de vrais et de faux positifs dans $r$, alors

  $$\text{précision} = \frac{t_p(r)}{t_p(r) + f_p(r)} = \frac{t_p(r)}{|r|}$$

  \vfill

  Le \alert{rappel} mesure la fraction de faux négatifs. 

  \medskip

  $$\text{rappel} = \frac{t_p}{t_p + f_n}$$

  \vfill

  L'évaluation d'un système de RI est difficile : implique des tests rigoureux
  avec des utilisateurs sur un échantillon.

\end{frame}
\end{comment}

\begin{frame}{Faux-positifs et négatifs, précision et rappel}
  \begin{center}
    \includegraphics[height=7cm]{fauxpositifs.png}
  \end{center}
\end{frame}


\begin{comment}
\begin{frame}{Performances}
  On peut améliorer les performances des moteurs de recherche à l'aide de
  plusieurs techniques : 
  \begin{itemize}
    \item des requêtes plus structurées 
      \begin{itemize}
        \item requêtes booléennes, 
        \item expressions rationnelles, 
        \item proximité, 
        \item recherche d'expression
      \end{itemize}
    \item des résultats classés
      \begin{itemize}
        \item modèle vectoriel 
        \item PageRank
        \item analyse sémantique latente
      \end{itemize}
  \end{itemize}
\end{frame}

\end{comment}

\section{Index (ou liste) inversés}

\subsection{Modèle}

\begin{frame}{Exemple de base}

  Un ensemble (modeste) de documents nous servira de guide.

  \vspace{.5cm}

  \begin{tabular}{>{\color{structure.fg}}ll}
    $d_1$&Le loup est dans la bergerie.\\
    $d_2$&Le loup et le trois petits cochons.\\
    $d_3$&Les moutons sont dans la bergerie.\\
    $d_4$&Spider Cochon, Spider Cochon, il peut marcher au plafond.\\
    $d_5$&Un loup a mangé un mouton, les autres loups sont restés dans la bergerie.\\
    $d_6$&Il y a trois moutons dans le pré, et un mouton dans la gueule du loup.\\
    $d_7$&Le cochon est à 12 Euros le Kg, le mouton à 10 Euros/Kg.\\
    $d_8$&Les trois petits loups et le grand méchant cochon.\\
  \end{tabular}

\end{frame}

\begin{frame}{Le besoin et la solution}

  \begin{defblock}{Besoin}
    On veut chercher tous les documents parlant de loups, de moutons mais pas de
    bergerie.
  \end{defblock}

  \medskip

  \alert{Parcourir tous les documents ?} (\texttt{grep})
  \begin{itemize}
    \item potentiellement long;
    \item critère "pas de bergerie" n'est pas facile à traiter;
    \item autres types de recherche ('loup' doit être \alert{près} de 'mouton'") sont
      difficiles;
    \item comment \alert{classer} par pertinence les documents trouvés ?
  \end{itemize}

  \vfill
\begin{blocgauche}
  Structure spécialisée : la \alert{matrice d'incidence} et surtout son inversion.
\end{blocgauche}
\end{frame}

\begin{frame}{Matrice avec documents en ligne}

  On sélectionne un ensemble de mots (ou \alert{termes}), constituant notre
  \alert{vocabulaire}

  \vfill

  \figSlideWithSize{matrice-incidence1}{10cm}


  \vfill
\begin{blocgauche}
  Documents en ligne, termes en colonnes. Dans chaque cellule : 1 si le terme est dans
  le document, 0 sinon.
\end{blocgauche}
\end{frame}


\begin{frame}{Recherche (loup, mouton, et pas bergerie)}

On prend les vecteurs binaires des \alert{termes} (les colonnes).

\begin{itemize}
  \item Loup :      11001101
  \item Mouton :    00101110
  \item Bergerie :    01010011
\end{itemize}

\vfill

Puis :

\begin{itemize}
  \item ET logique sur les vecteurs de Loup et Mouton, on obtient  00001100.
  \item ET logique  avec le \alert{complément} du vecteur de Bergerie (01010111)
\end{itemize}

\vfill
\begin{blocgauche}
On obtient 00000100, d'où on déduit que la réponse est limitée au document $d_6$.

Opération binaires \alert{très efficace}, mais...
\end{blocgauche}
\end{frame}


\begin{frame}{Passons à grande échelle}

\begin{itemize}
 \item Un million de documents, mille mots chacun en moyenne.
 \item Disons 6 octets par mot, soit 6 Go (ce n'est pas une si grosse base que
   cela)
 \item Disons 500\,000 termes \alert{distincts}
\end{itemize}

% 
\smallskip

$\Rightarrow$ la matrice d'incidence a :
\begin{itemize}
  \item $10^6$ lignes et 500 000 colonnes soit $500 \times 10^9$ bits
  \item soit 62 Go approximativement
\end{itemize}

\vfill

Ne tient pas en mémoire, ce qui va beaucoup compliquer les choses....

\vfill

\alert{Comment faire mieux\,?}

\end{frame}

\begin{frame}{On peut faire mieux}

  \begin{minipage}[c]{.46\linewidth}
       Il vaut mieux avoir les termes en ligne pour disposer des vecteurs dans une zone
mémoire contigue.

\vspace*{1cm}

On parle de matrice inversée, et de \alert{liste inversée}.
Une liste par terme; dans chaque liste, 1 pour les documents contenant le terme.
   \end{minipage} \hfill
   \begin{minipage}[c]{.46\linewidth}
\figSlideWithSize{inverted-index-bool}{4.5cm}
   \end{minipage}
   
\begin{blocgauche}
\begin{defblock}{Important}
Mises à jour beaucoup plus difficiles car elles impliquent la réorganisation
d'une liste compacte. Prix à payer pour l'efficacité en \alert{lecture} (recherche).
\end{defblock}
\end{blocgauche}
\end{frame}


\begin{frame}{On peut encore faire mieux}

La matrice est \alert{creuse} : il n'y a que $10^9$ positions avec des 1, soit un sur 500.

\vfill


\figSlideWithSize{inverted-index}{6cm}

\vfill

\begin{blocgauche}
\begin{defblock}{Essentiel}
On place dans les cellules \alert{l'identifiant} du document. De plus chaque liste est
\alert{triée} sur l'identifiant du document.
\end{defblock}

\end{blocgauche}
\end{frame}

\begin{frame}{Index inversé}

La structure utilisée dans \alert{tous}  les moteurs de recherche.

\begin{itemize}
 \item Un \alert{répertoire} contient tous les \alert{termes}.
 \item Une \alert{liste} (inversée) est associée à chaque \alert{terme}, \alert{triée} par docId.
 \item  Chaque élément de la liste est appelé une \alert{entrée}.
\end{itemize}

\vfill

%\alert{Concept}\,: la notion de \alert{terme} (\alert{token} en anglais) est différente de celle de ``mot''. 
%\vfill

%\alert{Vocabulaire}\,: le répertoire est parfois appelé \emph{dictionnaire}\,; les listes sont des \emph{posting list} en anglais;
%les entrées sont des \emph{postings}.


\begin{blocgauche}
\alert{Efficacité}\,: le répertoire devrait toujours être en mémoire\,; les listes, autant que possible en mémoire, sinon
fichiers contigus sur le disque.
\end{blocgauche}
\end{frame}

\subsection{Opération de recherche avec liste inversée}

\begin{frame}{Documents parlant de \alert{loup} {\bf ET} de \alert{mouton}}


Algorithme par \alert{fusion} : on parcourt \alert{séquentiellement} les deux listes.

\vfill

On compare à chaque étape les docId\,; s'ils sont égaux\,: \alert{trouvé}!

\vfill

On avance sur la liste du plus petit docId.

\vfill

\only<1>{
\figSlideWithSize{brutus-et-cesar}{7.5cm}
}
\only<2>{
\figSlideWithSize{brutus-et-cesar-2}{7.5cm}
}
\only<3>{
\figSlideWithSize{brutus-et-cesar-3}{7.5cm}
}
\only<4>{
\figSlideWithSize{brutus-et-cesar-4}{7.5cm}
}
\only<5>{
\figSlideWithSize{brutus-et-cesar-5}{7.5cm}
}
\only<6>{
\figSlideWithSize{brutus-et-cesar-6}{7.5cm}
}
\only<7>{
\figSlideWithSize{brutus-et-cesar-7}{7.5cm}
}
\only<8>{
\figSlideWithSize{brutus-et-cesar-8}{7.5cm}
}
\only<9>{
\figSlideWithSize{brutus-et-cesar-9}{7.5cm}
}

\pause
%\vfill

%\alert{Un seul parcours suffit}\,: recherche linéaire, parcours séquentiel. Pas mieux.

\vfill

\begin{blocgauche}
C'est une recherche dite \alert{Booléenne}\,: pas de classement, résultat exact. 
On peut faire mieux...
\end{blocgauche}
\end{frame}

\begin{comment}
\begin{frame}[fragile]{Algorithme}
\begin{minted}[frame=lines,linenos=true,
               fontsize=\footnotesize,
               framesep=3mm]{php}
// Fusion de deux listes l1 et l2
function Intersect($l1, $l2)
{
  $résultat = [];
  // Début de la fusion des listes
  while ($l1 != null and $l2 != null) {
    if ($l1.docId == $l2.docId) {
      // On a trouvé un document contenant les deux termes
      $résultat += $l1.docId;
      // Avançons sur les deux listes
      $l1 = $l1.next; $l2 = $l2.next;
    }
    else if ($l1.docId < $l2.docId) {
      // Avançons sur l1
      $l1 = $l1.next;
    }
    else {
       // Avançons sur l2
      $l2 = $l2.next;
    }
  }
}
\end{minted}

\end{frame}

\begin{frame}[fragile]{Optimisation}
  \begin{itemize}
    \item Si l'on veut effectuer la recherche : 
\begin{minted}[frame=lines,
               baselinestretch=.9,
               fontsize=\footnotesize,
               framesep=2mm]{php}
               A AND B AND C
\end{minted}
\item Quelle manière de procéder va optimiser le résultat ?
  \pause
  \item on stocke la taille des listes et l'on commence par faire l'intersection
    des plus petites
  \medskip
  \item de nombreuses autres questions d'optimisation existent 
    \item exemple : stocker l'une des listes en mémoire et calculer les
      intersections à la volée, en lisant depuis le disque
  \end{itemize}
\end{frame}
\end{comment}

\begin{frame}{Premier bilan}

Fondements techniques : index inversé, tri et compaction, parcours linéaire.

\begin{itemize}
 \item Permet d'effectuer des recherches \alert{puissantes}
     et \alert{flexibles}.
 \item \alert{Garantissent une très grande efficacité.}
\end{itemize}
\vfill

\begin{description}
 \item[\alert{Quels termes}\,?] Quels termes indexe-t-on\,? Beaucoup moins facile que ça n'en a l'air...
 \item[\alert{Quelles requêtes}\,?] De la plus simple (``sac de mots'') à plus 
                 structurée 
%(index structuré, requêtes Boolénnes, recherche de phrases).
   \item[Performance.] compression, distribution, optimisation des accès,  mises à jour, etc.
  \item[\alert{Classement}\,?]  comment classer en premier les documents pertinents\,?
\end{description}

\vfill
\begin{blocgauche}
La suite dans les prochaines sessions
\end{blocgauche}

\end{frame}

